% Section 8: Threats to Validity
% Source: context/Discussion_and_Conclusions.md §D9
%
% EM-DASH PROHIBITION: Do not use --- or \textemdash. Use colons or semicolons.

\section{Threats to Validity}\label{sec:threats}

\textbf{Construct validity.}
The binary \DFR{} metric does not capture corruption severity, but RSA
verification is itself binary (any distance $>0$ is fatal), making the metric
appropriate. Synthetic JWTs may under-represent real-world complexity; the eco
arm mitigates this by running with the exact production schema
($\sim$1{,}080~chars, RS256 Azure AD B2C) and obtaining
\DFR{}~=~19.9\%, statistically indistinguishable from the synthetic result.

\textbf{Internal validity.}
Experiments used direct provider API calls; identical corruption patterns
across six independent provider SDKs make a middleware artifact implausible.
All models were accessed under pinned version identifiers during a fixed
collection window of 3--7~July~2026; any residual within-study drift is
accounted for by the GEE clustering structure. Provider back-end updates
between study start and end could introduce uncontrolled variance; the fixed
collection window and version pinning limit but do not eliminate this risk.
The tokenisation hypothesis is the most parsimonious explanation for the
observed length-gating pattern; alternative mechanisms (tool-call
serialisation artefacts, provider SDK transforms, safety-layer rewrites)
have not been experimentally ruled out. Causal attribution is out of scope
for this study, whose primary contribution is demonstrating that the problem
exists at meaningful scale and that an architectural bypass eliminates it.

\textbf{External validity.}
The 16 models cover eight tiers and six providers but do not represent
the full model landscape; new models may exhibit different failure profiles.
The eco arm's $\rho\!=\!0.740$ rank correlation supports transfer to at least one
production deployment; generalisation to other credential types or schemas
requires future work (Section~\ref{sec:discussion}).
Per-model eco breakdowns are in the replication package.

\textbf{Conclusion validity.}
Multiple comparisons across 16 models inflate Type~I error; all pairwise
comparisons are corrected by Benjamini-Hochberg at FDR~=~0.05. The primary
inference (\DFR{} $> 0$\% across all models) is a single test with effectively
unlimited power at $N\!=\!50{,}400$. The analysis plan was pre-registered
before data collection to reduce observer bias.
